% Figure: Sutton-Graves stagnation heat flux versus nose radius for two
% entry speeds, showing the 1/sqrt(R_n) falloff that makes bluntness the
% design lever. q = K sqrt(rho/Rn) U^3, K=1.83e-4 SI, rho=2e-5 kg/m3.
% U=7 km/s and U=11 km/s. q in MW/m^2 (divide W/m2 by 1e6).
% prefactor A = K*sqrt(rho)*U^3 / 1e6:
%  U=7e3:  1.83e-4*sqrt(2e-5)*(7e3)^3/1e6 = 1.83e-4*4.4721e-3*3.43e11/1e6
%          = 1.83e-4*4.4721e-3=8.184e-7; *3.43e11=280.7; /1e6=2.807e-4? recompute below in comment
% Careful: 8.184e-7*3.43e11 = 2.807e5 W/m2 base (at Rn=1). /1e6 = 0.2807 MW/m2 at Rn=1.
%  so plot A/sqrt(Rn): U=7 -> 0.281/sqrt(Rn); U=11 -> scale by (11/7)^3=3.88 -> 1.090/sqrt(Rn)
\begin{figure}[htbp]
\centering
\begin{tikzpicture}
\begin{axis}[
  width=12cm, height=7.5cm,
  xlabel={nose radius $R_n$ (m)},
  ylabel={stagnation heat flux $q_w$ (MW/m$^2$)},
  xmin=0.05, xmax=2.0, ymin=0, ymax=5,
  axis lines=left,
  grid=major, grid style={enggray!20},
  tick label style={font=\scriptsize},
  label style={font=\small},
  legend style={font=\scriptsize, at={(0.97,0.97)}, anchor=north east},
  legend cell align=left,
]
% U = 11 km/s (lunar-return class)
\addplot[engred, very thick, domain=0.05:2.0, samples=120] { 1.090/sqrt(x) };
\addlegendentry{$U_\infty=11\,$km/s (lunar return)}
% U = 7 km/s (low-Earth-orbit return)
\addplot[engblue, very thick, domain=0.05:2.0, samples=120] { 0.281/sqrt(x) };
\addlegendentry{$U_\infty=7\,$km/s (LEO return)}
% mark a representative shuttle leading-edge radius
\draw[enggray, dashed] (axis cs:0.3,0) -- (axis cs:0.3,5);
\node[enggray, font=\scriptsize, anchor=south west, rotate=90] at (axis cs:0.3,0.2)
  {$R_n\approx0.3$ m (leading edge)};
\end{axis}
\end{tikzpicture}
\caption[Stagnation heat flux versus nose radius]{Stagnation-point heat
flux from the Sutton-Graves correlation $q_w=K\sqrt{\rho_\infty/R_n}\,
U_\infty^3$ at a peak-heating density $\rho_\infty\approx2\times10^{-5}\,
\si{\kilogram\per\meter\cubed}$, for a low-Earth-orbit return at
$7\,\si{\kilo\meter\per\second}$ (blue) and a lunar return at
$11\,\si{\kilo\meter\per\second}$ (red). The flux falls as
$1/\sqrt{R_n}$, so doubling the nose radius cuts the peak flux by about
thirty percent; this inverse-square-root dependence is the quantitative
reason re-entry vehicles are blunt. The cube on velocity is the steeper
scaling: the lunar-return curve sits a factor $(11/7)^3\approx3.9$ above
the orbital-return curve at every radius, which is why a return from the
Moon or a planet is a far harder thermal problem than a return from low
orbit. The dashed line marks a reinforced-carbon-carbon leading-edge
radius near $0.3\,\si{\meter}$.}
\label{fig:vol03ch13:heat-flux-radius}
\end{figure}
